\section{Introducción}

Los mercados de criptoactivos son poco regulados y tienen baja liquidez, lo
que los hace vulnerables a manipulación. Uno de los esquemas más comunes es el
\emph{pump \& dump}: un grupo coordinado compra masivamente un token para
inflar su precio y luego lo vende, dejando pérdidas a los inversores que
entraron tarde~\cite{xu2019anatomy}. Se estima que el 24\,\% de los tokens
lanzados en 2022 tuvo características de este tipo de
manipulación~\cite{lamorgia2023}.

Los métodos actuales de detección analizan cada token por separado: buscan
picos anómalos de precio o volumen~\cite{kamps2018}, o usan modelos de
aprendizaje automático~\cite{hamrick2021}. El problema es que la manipulación
coordinada deja una huella colectiva que estos enfoques no pueden ver: varios
tokens que de repente se mueven juntos de forma artificial, formando una
comunidad densa en la red de correlaciones.

Este trabajo propone detectar esa huella usando análisis espectral de grafos.
Representamos el mercado como una red donde los nodos son tokens y las aristas
reflejan qué tan correlacionados están sus precios. Luego analizamos la
estructura matemática de esa red para encontrar agrupaciones artificiales.

La herramienta central es el \emph{Inverse Participation Ratio} (IPR),
un concepto originado en física~\cite{anderson1958absence} que mide qué tan
concentrada está la energía en un modo de vibración de la red. Trabajos
previos lo usan de forma descriptiva sobre matrices de
correlación~\cite{plerou1999,nobi2013}; este trabajo propone usarlo por primera
vez como detector de manipulación, con validación sobre eventos reales
etiquetados.
